% 本文件由 LaTeX 排版 Agent 根据有效 translation.json 自动生成。
% author=Tertullian; work_id=0314; title=驳瓦伦提尼安派

\chapter{驳瓦伦提尼安派}

% structure: p0001 source=h1 target=omitted reason=page-title-already-rendered-by-parent

% structure: p0002 source=h2 target=section reason=relative-heading-rank; base=h2

\section{第一章 导言：戴尔都良将此异端比作古老的厄琉息斯秘仪。二者同样宁愿隐藏错误和罪恶，而不愿宣讲真理与德行。}

瓦伦提尼安派无疑是一大批异端者——他们由众多背弃真理者组成，倾向寓言，且无纪律能阻止他们——他们最关心的莫过于模糊自己所传讲的，如果那确实可称为传讲的话，因为模糊\Emph{他们的教义}的人就是在传讲。他们守护教义的殷勤正是一种显露罪过的殷勤。他们维护其宗教体系的极大热诚恰恰宣扬了他们的耻辱。至于那些厄琉息斯秘仪，正是雅典迷信的异端，其耻辱就在于它的隐秘。因此，他们预先用折磨人的条件封锁进入他们团体的一切途径，\Emph{并且}要求长期入门，甚至为成全的门徒进行五年的教导，以便通过这种对完全知识的延宕来塑造他们的意见，并显然因预先制造了对秘仪的渴望而提升其尊严。随之而来的是沉默的义务。那长久寻觅的东西被谨慎守护。然而，一切神性都在他们的秘密内室里：在那里\Emph{最终}启示出已完全入门者的所有渴望，密封之舌的全部奥秘，以及男子气概的象征。但这种寓意描绘，借自然之尊严名义的托辞，用任意的象征掩盖真正的渎圣，并借\Emph{空洞}形象避免谎言的谴责！同样，目前我们讨论的这些异端者，即瓦伦提尼安派，已形成了他们自己的厄琉息斯式荒宴，以深深的沉默圣化，除了神秘之外毫无天上的事物。借助真宗教的神圣名称、名号和论证，他们为人们柔顺的喜好编造出最虚妄最污秽的虚构，出自圣经的丰富暗示，因为从许多泉源中可能涌出许多\Emph{错误}。你若向他们提出真诚正直的询问，他们便以严峻的目光和皱起的眉头回答并说：\Quote{课题深奥。}你若以微妙的问题、用他们双舌的含糊语试探他们，他们便肯定与你有共同的信德。你若暗示你理解\Emph{他们的意见}，他们便坚持说自己一无所知。你若与他们短兵相接，他们便以自我献祭摧毁你自己对胜利的盲目希望。他们甚至在确保门徒忠诚之前，不将秘密托付给门徒。他们具有在教导之前先说服人的技巧，尽管真理通过教导说服人，却不先说服人而后教导。\par

% structure: p0004 source=h2 target=section reason=relative-heading-rank; base=h2

\section{第二章 这些异端者给基督徒贴上\Quote{单纯者}的标签。接受此指控，并基于圣经颂扬单纯。}

为此，他们给我们贴上了单纯者的标签，而且只是单纯，并无智慧；仿佛智慧必须缺少单纯，而主却将两者结合：\BibleQuote{所以你们要机警如同蛇，纯朴如同鸽子。}\BibleRef{玛 10:16} 若我们因单纯而被视为愚拙，那么他们因有智慧是否就不单纯呢？然而，最乖谬的正是那些不单纯者，正如最愚拙的正是那些不智慧者。但（若要我选择），我宁愿接受\Emph{后一种}状态为较小的过失；因为或许有数量不足的智慧，比品质恶劣的智慧更好——犯错比引人误入歧途更好。此外，那些\BibleQuote{以纯朴的心寻求祂}的人耐心等候主的慈颜，这是那智慧之言——不是瓦伦提诺的，而是撒罗满的——所说的。\BibleRef{智 1:1} 再者，婴儿曾以自己的血为基督作证。难道（你说）是孩童呼喊\Quote{钉死祂}吗？他们既非孩童亦非婴儿；换言之，他们并不单纯。宗徒也吩咐我们要向着天主\Quote{重新成为婴孩}，因我们的单纯而\Quote{在邪恶上成为婴孩}，但也要在实践能力上\Quote{成为智慧}。同时，关于智慧的发展次序，我承认它源于单纯。简言之，\Quote{鸽子}通常用以预表基督；\Quote{蛇}则用以诱惑祂。前者从起初就是神圣和平的使者；后者从起初就是神圣形象的抢夺者。因此，单单单纯将更能认识并宣告天主，\Emph{而}单单智慧则会行暴于祂并背叛祂。\par

% structure: p0006 source=h2 target=section reason=relative-heading-rank; base=h2

\section{第三章 此异端的愚妄。它分割并肢解天主。与真宗教的单纯智慧对比。揭露瓦伦提尼安体系的荒谬就是摧毁它。}

让蛇尽其所能地隐藏自己，将其全部智慧扭曲在谜团的迷宫中；让它深居地下，钻进隐秘的洞穴；让它通过蜿蜒的关节展开身躯；让它曲折地爬行，尽管并非一次完成，这畏光的野兽。然而，我们的鸽子，它的巢穴是何等单纯！——总是在高处和开阔之处，面向光明！作为圣神的象征，它喜爱那（光明的）东方，那是基督的预像。真理毫无可羞之处，除非被隐藏，因为没有人会以倾听真理为耻。\Emph{没有人会以承认祂为天主为耻}，因为本性已经向他们推荐了祂，他们在祂的一切工程中已经感知了祂——确实，祂之所以不为人完全认识，仅仅是因为人没有认为祂是唯一的一位，因为他以复数（众神）称呼祂，并以其他\Emph{形态}朝拜祂。然而，让自己从这一大堆神祇转向另一群人，从熟悉的权威移向未知的权威，从明显之事转向隐藏之事，这是在信仰的门槛上冒犯信仰。现在，即使你完全接受了那一整套神话，难道你不会想起，在你婴儿时期，你慈爱的保姆在你耳边唱摇篮曲时，给你讲过类似的故事，关于拉弥亚的塔楼和太阳的角吗？然而，让任何人从他所学到的信仰知识来接触这个问题，一旦他发现有如此多的永世的名称，如此多的婚姻，如此多的后裔，如此多的出口，如此多的结局，以及一个分散和残缺的天主的福\Emph{与}祸，他难道会不立即断定，这些就是蒙受默感的宗徒预先谴责的\Quote{神话和无穷的族谱}，而这些异端的种子那时就已经萌发了吗？因此，应当公正地认为，那些制造这些神话并非毫不费力、只间接为之辩护、同时又不彻底教导他们所教之人的人，缺乏单纯而只属精明。当然，这显示了他们的狡猾，如果他们的教训是可耻的；也显示了他们的刻薄，如果这些教训是光荣的。至于我们这些单纯的人，我们完全了解这一切。简言之，这是我们为交锋所装备的第一件武器；它揭露并展示他们整个败坏的系统。这是我们胜利的第一个预兆；因为仅仅指出那以如此大的技巧隐瞒的东西，就是摧毁它。\par

% structure: p0008 source=h2 target=section reason=relative-heading-rank; base=h2

\section{第四章。异端可追溯到瓦伦廷，一个能干但不安分的人。该学派的许多分裂领袖被提及。只有一人尊重以其名字命名整个学派的人}

我说，我们完全知道他们真正的起源，也清楚为什么称他们为瓦伦廷派，尽管他们假装否认自己的名字。他们确实离开了他们的创始人，但他们的起源决没有被消灭；即使它偶然改变了，这改变本身也证明了事实。瓦伦廷曾期望成为主教，因为他是一个在天才和口才上都颇有能力的人。然而，因着另外一人因宣认信仰所赋予的权利而获得这尊严，他愤而脱离了真信仰的教会。就像那些（不安的）精神，被野心激发时通常被复仇的欲望点燃，他全力扑灭真理；并找到某种古老见解的线索，他以蛇的诡计为自己开辟了一条道路。后来，托勒密走上了同一条道路，将永世的名字和数目区分为位格性的实体，但他却使这些实体与天主分开。瓦伦廷曾将这些实体包含在天主本质之内，作为情感的知觉和感动。然后，由赫拉克林、塞昆杜斯和术士马可从中开辟了各种旁路。特奥提摩斯致力于\Quote{法律的预像}。然而，瓦伦廷本人尚无处可寻，而瓦伦廷派仍从瓦伦廷得名。安提约基雅的阿克修尼库斯是当今唯一充分保持瓦伦廷规则、以尊崇其记忆的人。但这个异端被允许像娼妓一样以各种形态装扮自己，娼妓通常每天更换和调整她的服饰。为何不呢？当他们以这种方式审视每个人里面的那属灵的种子时，一旦他们想到什么新奇事物，就立即称他们的妄称为启示，他们自己的乖戾巧智为神恩；但他们（否认一切）的统一，只\Emph{承认}多样性。因此我们清楚看到，撇开他们惯常的伪饰，他们大多数人处于分裂状态，随时准备（且真诚地）对他们的信仰的某些观点说：\Quote{不是这样}；\Quote{我以不同方式理解}；\Quote{我不承认这一点}。确实，通过这种多样性，创新就在他们的规则的表面上打上了印记；此外，它还披着无知的自负的一切貌似合理的色彩。\par

% structure: p0010 source=h2 target=section reason=relative-heading-rank; base=h2

\section{第五章。许多杰出的基督徒作家已经仔细而充分地驳斥了这异端。作者将这些作家作为自己的向导}

然而我的道路是沿着他们首席导师的最初教义，而非他们混杂跟随者中的自封领袖。我们不会从任何方面听到人说，我们凭己意捏造了材料，因为关于这些意见及其反驳，已有许多极其圣善卓越的人——不仅在我们之前生活的人，也包括与异端首领同时代的人——精心写成著作予以呈现：例如殉道者哲学家\Person{犹斯定}、教会的诡辩家\Person{米耳提雅德}、对一切学说探究极为精确的\Person{依勒内}、我们自己的\Person{普罗库鲁斯}——贞洁老年与基督宗教雄辩的典范。我渴望在每一项信仰工作中都紧紧跟随他们，正如在这一特别工作中一样。如果根本没有异端，而那些反驳异端的人却被认为捏造了它们，那么预言这些异端的宗徒\BibleRef{格前 11:19}就必犯了虚谎之罪。然而，若确有异端，则只能是这些正在讨论中的异端。不能认为任何作家有如此多的时间去捏造已经在他掌握中的材料。\par

% structure: p0012 source=h2 target=section reason=relative-heading-rank; base=h2

\section{第六章 虽然以拉丁文写作，他提议保留华伦提努派神圣流溢的希腊文名称。并不讨论异端，只揭露它。对此加以其荒谬应得的嘲弄。}

那么，为使无人被如此众多奇特的名字所迷惑——这些名字被收集起来，随意调整，含义不明——我在这篇短文中仅仅承担阐述这个（异端）奥秘的任务，说明我们应当如何使用它们。现在，将其中一些\Emph{名字}从希腊文翻译过来以产生同样明显的词义，绝非易事：有些名字的性不相配；另一些则以希腊文形式更为人所熟知。因此，大部分时候我们将使用希腊文名称；其含义将见于页边空白。希腊文也并非没有拉丁文的\Emph{对应词}；只是这些对应词将标记在上方行间，以说明人名——这是必要的，因为其中有些名字含混，容许不同的意义。但是，虽然我必须搁置一切讨论，目前只满足于（对异端的）单纯阐述，但无论何处若有丑闻特征似乎需要惩戒，就须以一切方式攻击它，即使只是轻刺一下也罢。让读者视之为战前的小冲突。我的意图是显示如何刺伤而非造成深创伤。若某处引发了欢笑，这恰恰是主题应得的。有许多事物需要以不耗费庄重的方式加以驳斥。虚荣愚蠢的话题特别适合以笑声应对。甚至真理也可以纵情嘲笑，因为它喜乐；它可以与敌人戏耍，因为它无畏。只是我们必须当心，它的笑声不可失当，以致自身被嘲笑；但凡其欢笑得体，便是一种义务去\Emph{纵容它}。如此，我终于开始我的任务。\par

% structure: p0014 source=h2 target=section reason=relative-heading-rank; base=h2

\section{第七章 最初的八种流溢，或称永恒存在，称为八元体，是其余一切之源。其名号与谱系记录。}

从罗马诗人\Person{恩纽斯}开始，他仅仅言及\Quote{天堂的宽敞大厅}——或是由于它们位置高耸，或是因为在\Person{荷马}处他读到\Person{朱庇特}在其中宴饮。然而，至于\Emph{我们的}异端者，他们悬挂、升起并展开一层又一层、一重又一重的高度，作为他们各自神祇的居所，这真是奇妙。甚至我们的造物主也已被安排了恩纽斯式的大厅，如同私人房间，房间叠着房间，并且按照异端的数目给每位神祇分配了同样多的阶梯。\Emph{事实上}，宇宙已变成了\Quote{出租的房间}。你会想象这些天堂的楼层是某处不知名的福乐岛上独立的寓所。在那里，甚至华伦提努派的神也住在阁楼里。他们确实按其实质称他为\OriginalTerm{完美永恒}{\GreekText{Αἰῶν} \GreekText{τέλειος}}（\Emph{完美永恒}），但就其位格而言，则称他为\OriginalTerm{太初}{\GreekText{Προαρχή}}（\Emph{太初}）。\par

\SourceRecord{Translated by Alexander Roberts.；Ante-Nicene Fathers；Vol. 3.；Edited by Alexander Roberts, James Donaldson, and A. Cleveland Coxe.；Buffalo, NY: Christian Literature Publishing Co.,；1885.；<http://www.newadvent.org/fathers/0314.htm>.}
